% LSH 核心术语定义
% 版本: v5.0
% 更新日期: 2026-05-11

% ============================================
% 第一部分：基础概念
% ============================================

\def\LSH{LSH (Li Shi Hang)}
\def\LSHSFA{LSH-SFA (Spacetime Field Attention)}
\def\LSHID#1{\texttt{#1}}

% ============================================
% 第二部分：核心表示
% ============================================

\def\WavePacket{波包表示 (Wave Packet Representation)}
\def\WavePacketShort{波包}
\def\LightCone{光锥注意力 (Light-Cone Attention)}
\def\FieldAttention{场注意力 (Field Attention)}
\def\PDEEvo{PDE 驱动的场演化 (PDE-Driven Field Evolution)}
\def\ResonanceCoupling{共振耦合 (Resonance Coupling)}

% ============================================
% 第三部分：三层架构
% ============================================

\def\ThreeLayer{三层语义架构 (Three-Layer Architecture)}
\def\StructureLayer{结构层 (Structure Layer)}
\def\RuleLayer{规则层 (Rule Layer)}
\def\ExecutionLayer{执行层 (Execution Layer)}
\def\Element{元素 (Element)}
\def\Observer{观察者 (Observer)}
\def\CMDPerceive{CMD\_PERCEIVE}

% ============================================
% 第四部分：梯度系统
% ============================================

\def\VTwo{V2 无状态梯度架构 (V2 Stateless Gradient)}
\def\GradientParadigm{梯度范式研究 (Gradient Paradigm)}
\def\ControllableExplosion{梯度爆炸 (Gradient Explosion)}
\def\EnergyConservation{能量守恒 (Energy Conservation)}

% ============================================
% 第五部分：时空场
% ============================================

\def\SpacetimeEngine{时空场引擎 (Spacetime Engine)}
\def\DiffusionEq{扩散方程 (Diffusion Equation)}
\def\ConvectionEq{对流方程 (Convection Equation)}
\def\WaveEq{波动方程 (Wave Equation)}

% ============================================
% 使用示例
% ============================================
%
% 在论文中使用：
%   \LSH 提出了...
%   \WavePacket 由四元组 (k, w, \omega, \theta) 组成
%   \ThreeLayer 包括 \StructureLayer、\RuleLayer 和 \ExecutionLayer
%
